%-----------------------------------------------------------------------
%  Packages, column type and the \hd heading macro.
%  \input by main.tex -- not compiled on its own.
%  Shares no file with ../../master_plan/ or with any other brief, so this
%  document can be its own Overleaf project and go to one reader alone.
%-----------------------------------------------------------------------
\usepackage[T1]{fontenc}
\usepackage[utf8]{inputenc}
\usepackage{lmodern}
\usepackage[margin=0.85in]{geometry}
\usepackage{amsmath}
\usepackage{booktabs}
\usepackage{tabularx}
\usepackage{enumitem}
\usepackage{microtype}
\usepackage{graphicx}

\newcolumntype{Y}[1]{>{\hsize=#1\hsize\raggedright\arraybackslash}X}
% \par first: \hd is a run-in heading and must open its own paragraph. Without
% it a heading at the top of an \input file joins the previous file's last
% paragraph, which is invisible in the source and obvious in the PDF.
\newcommand{\hd}[1]{\par\smallskip\noindent\textbf{#1}\ \ignorespaces}
\newcommand{\code}[1]{\texttt{\small #1}}
\newcommand{\kv}{\ensuremath{\kappa_{v}}}
% \dg is \ensuremath, not $^{\circ}$: that form is text-mode only, so \dg
% inside math closes the surrounding math and reopens it, and the error points at
% the end of the enclosing table rather than at the offending cell. Renders
% identically in text mode.
\newcommand{\dg}{\ensuremath{^{\circ}}}

% Float pages (the full-page synthesis figure) anchor to the top of the page
% rather than centring vertically, which reads better with a long caption.
\makeatletter
\setlength{\@fptop}{0pt}
\makeatother

\setlength{\parindent}{0pt}
\setlength{\parskip}{1.5pt}
\renewcommand{\arraystretch}{1.06}
\pagestyle{empty}
